\section{Public Artifact Stack And Lossy Projection}

Commitments and public shadows are the first concrete public-facing objects in
the paper. They explain what public data can carry without revealing the hidden
Hodge representative.
They are the scaffolding for transparency in action and anonymity in identity:
actions are exposed through auditable public artifacts, while identities remain
inside hidden representatives and public equivalence classes.

\subsection{Public, Artifact, And Private Layers}

The public-facing stack has three mathematical levels. The fully public handle
level contains hashes, roots, topology hashes, nullifiers, public summaries, and
oracle judgments. The verifier-facing artifact level contains public shadows,
cochain commitments, proof-carrying cochain payloads, and certificate data. The
private witness level contains the fibration representative,
transport/restriction data, full cochains and openings, and the Hodge operator
whose kernel selects the harmonic component. The paper's claims are stated so
that public handles and verifier-facing artifacts bind to hidden witnesses
without turning the hidden witness into public data. This is compatible with the
sheaf-theoretic attacker-observation viewpoint of
\cite{inoue2026cryptographicSheaves}; the associated security games and
runtime observation model belong to the subsequent protocol paper. The present
paper uses only the mathematical public/private separation described below.

\subsection{The Artifact Stack}

The public artifact stack has the following shape:
\[
\begin{aligned}
  &\text{hidden/private object} \\
  &\quad\to \text{public shadow or compressed public datum} \\
  &\quad\to \text{commitment/root/hash artifact} \\
  &\quad\to \text{proof-carrying artifact} \\
  &\quad\to \text{verifier judgment}.
\end{aligned}
\]
Each layer binds to or summarizes the preceding layer. A hash does not replace
the mathematical object, a commitment does not replace the witness, and a proof
artifact does not replace the proof obligation. The purpose of the stack is to
let public verification refer to hidden sheaf/cochain data without exposing it.
In this sense, the stack separates what must be publicly accountable from what
must remain representative-private.

Hashes serve as public names, integrity bindings, and fixed-size pointers for
larger artifacts. A topology hash names the public combinatorial support; a
root commits to a batch or storage view; a transcript hash domain-separates a
proof statement. None of these hashes asserts validity by itself. Validity
comes from the predicate or proof attached to the hashed object.

\subsection{Cochain Commitments}

Cochain commitments provide the public handle on hidden cochain data. The Lean
surface treats the canonical payload as a mathematical object with a modulus,
edge values, declared edge count, and topology hash.

\subsection{Commitment Map}

\begin{definition}[Cochain commitment payload]\label{def:cochain-commitment-payload}
  A cochain commitment payload is a tuple
  \[
    (q,n,c,\tau)
  \]
  consisting of a modulus \(q\), a declared edge count \(n\), a list \(c\) of
  integer edge values, and a topology hash \(\tau\). In the canonical Lean
  payload this is \texttt{CochainCommitmentData}.
\end{definition}

The intended commitment map is the transaction-facing coboundary map
\[
  (m,r) \longmapsto \delta(m\Vert r) \in C^1 \pmod q,
\]
stored together with the topology hash of the supporting complex. The payload
therefore preserves the accepted modulus, the edge-vector shape, the public
topology binding, and the serialized \(C^1\)-level commitment values. It does
not expose the hidden fibration representative, private transport data, the
operative Hodge--Laplacian, or the harmonic kernel.

\begin{definition}[Commitment well-formedness]\label{def:commitment-well-formedness}
  A cochain commitment payload is well formed when its modulus is the supported
  canonical modulus, its declared edge count agrees with the actual edge-value
  list length, and the edge list is nonempty.
\end{definition}

The security interface for this commitment layer consists of binding, hiding,
homomorphic addition, homomorphic scalar multiplication, topology binding, and
partial-opening soundness. The blinded-commitment, vector-commitment, and
cohomological-MAC interfaces are separate axiom records rather than consequences
silently bundled into the payload shape.

\subsection{Commitment And Artifact Invariants}

The commitment-facing invariants attach to the commitment payload and to the
public artifact stack, not to the proximity projection. Payload
well-formedness is a shape condition: it asserts the supported modulus,
agreement between the declared edge count and the edge-value list, and a
nonempty edge list. It is not a cryptographic uniqueness theorem and it does
not assert that the cochain vector is nonzero. For example, the canonical zero
payload is well formed whenever the declared edge count is positive.

Topology compatibility is a separate public invariant. Two committed cochains
can be aggregated only when their moduli, declared edge counts, and topology
hashes agree. The weaker same-topology relation records equality of topology
hashes alone. Cryptographic binding and topology-binding claims are supplied by
the commitment-security axioms, not by the raw payload record.

At the artifact level, Lean records a hashed committed cochain as a public
cochain payload together with digest data and a binding predicate. A
proof-carrying committed cochain adds a proof system, a statement, and a
verification predicate. The basic theorem exposed to this paper is that a
well-formed proof-carrying committed-cochain artifact has an underlying
well-formed cochain commitment payload.

\subsection{Sheaf Storage And Witness Material}

Storage roots, topology hashes, cochain roots, nullifiers, and proof objects are
public artifacts in the stack. Full sheaf, cochain, and fibration data remain
witness material unless a protocol path explicitly exposes them. The storage
mechanics are protocol-level concerns; the mathematical point here is that
stored public artifacts are shadows of accepted witness data, not replacements
for the hidden object.

\subsection{Lossy Projection}

Compression creates verifier-facing public artifacts from hidden or full
witness data. The compression need not be injective; it must preserve enough
structure to support the relevant public statement.

A public shadow is allowed to forget hidden data. In fact, it must forget
enough data that it does not select a private representative. The relevant
question is not whether the public datum is injective, but which facts survive
the compression relation.

\subsection{Preserved And Lost Structure}

The preserved structure before the finite-code layer is:
\begin{enumerate}[label=(\roman*)]
  \item public classifier or fingerprint equality for hidden representatives;
  \item cochain commitment payload well-formedness;
  \item topology binding for cochain commitments;
  \item digest/root binding for compressed public shadows;
  \item proof-carrying validity of the public statement attached to an artifact.
\end{enumerate}

The lost structure is the private representative itself: full sheaf/cochain
witness data, hidden transport, the operative Hodge operator, private
randomness, and runtime lifecycle facts. Those may be bound to public artifacts
or certified by a proof-carrying statement, but they are not made public merely
by the presence of a hash, root, or commitment payload.

\subsection{Fundamental Properties}

\begin{proposition}[Well-formed committed payloads]\label{prop:well-formed-committed-payloads}
  A well-formed cochain commitment payload has nonempty edge data and modulus
  greater than one.
\end{proposition}

\begin{proposition}[Topology-compatible aggregation]\label{prop:topology-compatible-aggregation}
  Commitment aggregation is only meaningful for payloads with compatible
  modulus, declared edge count, and topology hash. This is the public
  topology-binding condition consumed by the chain layer.
\end{proposition}

\begin{proposition}[Proof-carrying artifact exposes payload]\label{prop:proof-carrying-artifact-exposes-payload}
  A well-formed proof-carrying committed-cochain artifact exposes a well-formed
  underlying cochain commitment payload.
\end{proposition}

Byte parsing, allocation bounds, canonical encoding, concrete hash functions,
domain separation, storage inclusion proofs, proof verification, state updates,
and ephemeral-fiber lifecycle properties are TCB obligations. Lean records
their boundary predicates; it does not identify them with the mathematical
cochain commitment payload.
